\documentclass[12pt,a4paper]{report}

% ------------------------------------------------------------
% KÓDOLÁS, NYELV, BETŰKÉP
% ------------------------------------------------------------
\usepackage[T1]{fontenc}
\usepackage[utf8]{inputenc}
\usepackage[magyar]{babel}
\usepackage{lmodern}
\usepackage{microtype}

% ------------------------------------------------------------
% OLDALBEÁLLÍTÁS
% ------------------------------------------------------------
\usepackage[
a4paper,
left=25mm,
right=25mm,
top=25mm,
bottom=25mm
]{geometry}

% ------------------------------------------------------------
% MATEMATIKA ÉS MÉRTÉKEGYSÉGEK
% ------------------------------------------------------------
\usepackage{amsmath}
\usepackage{amssymb}
\usepackage{mathtools}
\setcounter{MaxMatrixCols}{20}
\usepackage{bm}
\usepackage{siunitx}
\sisetup{output-decimal-marker={,}}

% ------------------------------------------------------------
% TÁBLÁZATOK, LISTÁK
% ------------------------------------------------------------
\usepackage{booktabs}
\usepackage{longtable}
\usepackage{array}
\usepackage{enumitem}

% ------------------------------------------------------------
% ÁBRÁK ÉS BLOKKVÁZLATOK
% ------------------------------------------------------------
\usepackage{tikz}
\usetikzlibrary{arrows.meta,positioning,calc,shapes.geometric}
\usepackage{float}

% ------------------------------------------------------------
% KÓDRÉSZLETEK
% ------------------------------------------------------------
\usepackage{xcolor}
\usepackage{listings}
\lstdefinestyle{matlab}{
	language=Matlab,
	basicstyle=\ttfamily\small,
	keywordstyle=\bfseries,
	commentstyle=\itshape,
	showstringspaces=false,
	breaklines=true,
	frame=single,
	columns=fullflexible
}

% ------------------------------------------------------------
% HIVATKOZÁSOK
% ------------------------------------------------------------
\usepackage[hidelinks]{hyperref}

% ------------------------------------------------------------
% SAJÁT PARANCSOK
% ------------------------------------------------------------
\newcommand{\mbody}{m_{\mathrm{b}}}
\newcommand{\mwheel}{m_{\mathrm{w}}}
\newcommand{\Jspin}{J_{\mathrm{a}}}
\newcommand{\Jzw}{J_{z,\mathrm{w}}}
\newcommand{\tauL}{\tau_{\mathrm{L}}}
\newcommand{\tauR}{\tau_{\mathrm{R}}}
\newcommand{\tauc}{\tau_{\mathrm{c}}}
\newcommand{\taud}{\tau_{\mathrm{d}}}
\newcommand{\phiL}{\varphi_{\mathrm{L}}}
\newcommand{\phiR}{\varphi_{\mathrm{R}}}
\newcommand{\omegaL}{\omega_{\mathrm{L}}}
\newcommand{\omegaR}{\omega_{\mathrm{R}}}
\newcommand{\alphL}{\alpha_{\mathrm{L}}}
\newcommand{\alphR}{\alpha_{\mathrm{R}}}
\newcommand{\thetad}{\dot{\theta}}
\newcommand{\thetadd}{\ddot{\theta}}
\newcommand{\psid}{\dot{\psi}}
\newcommand{\psidd}{\ddot{\psi}}

\begin{document}
	
	\title{KÉTKEREKŰ ÖNKIEGYENSÚLYOZÓ ROBOT 3D DINAMIKAI MODELLEZÉSE ÉS IRÁNYÍTÁSA KLASSZIKUS ÉS TANULÁSALAPÚ SZABÁLYOZÁSI MÓDSZEREKKEL}
	
	\author{
		\textbf{Várkonyi Borisz}\\\\
		R83C3F \\
		Mechatronikai mérnöki alapképzési szak \\
		Nappali \\
		Szent István Campus
		\and
		\textbf{Dr. Erdélyi Viktor Ferenc}\\\\
		Egyetemi Adjunktus \\
		MATE, Mechatronika Tanszék\\
		Szent István Campus
	}
	\date{2026}
	
	\maketitle
	
	\chapter{Bevezetés}
	
	\chapter{Szakirodalom áttekintés}

	\chapter{Anyag és módszer}
	
	A fejezetben bemutatom a saját munkát
	
	\section{A vizsgált rendszer mechanikai modellje}

	A fejezet célja bemutatni a rendszer mechanikai modelljét.
	
	\subsubsection{Koordinátarendszer és előjel-konvenció}
	
	A világ koordinátarendszere:
	
	\begin{align*}
		X &: \text{előre},\\
		Y &: \text{balra},\\
		Z &: \text{felfelé}.
	\end{align*}
	
	
	A robot helyzetét és orientációját az alábbi mennyiségekkel írjuk le:
	
	\begin{itemize}
		\item $x$: a keréktengely középpontjának $X$ koordinátája;
		\item $y$: a keréktengely középpontjának $Y$ koordinátája;
		\item $\psi$: a robot irányszöge a függőleges tengely körül;
		\item $\theta$: a főtest hosszirányú dőlésszöge;
		\item $\phiL$: a bal kerék talajhoz viszonyított abszolút gördülési szöge;
		\item $\phiR$: a jobb kerék talajhoz viszonyított abszolút gördülési szöge.
	\end{itemize}
	
	Előjel-konvenció:
	
	\begin{itemize}[leftmargin=2em]
		\item $\theta>0$: a főtest teteje előrefelé dől;
		\item $\psi>0$: felülről nézve a robot balra fordul, vagyis pozitív $Z$ tengely körüli elfordulás történik.
	\end{itemize}
	
	A teljes konfigurációs vektor:
	
	\begin{equation}
		q_{teljes} = \begin{bmatrix}
			x & y & \psi & \theta & \phiL & \phiR
		\end{bmatrix}^{\!T}.
	\end{equation}
	
	\newpage
	
	\subsubsection{A modell fizikai paraméterei}
	
	\renewcommand{\arraystretch}{1.4}
	
	\begin{longtable}{>{\centering\arraybackslash}p{2.5cm} p{8cm} p{2.5cm}}
		\toprule
		\textbf{Jel} & \textbf{Jelentés} & \textbf{Egység} \\
		\midrule
		\endfirsthead
		\toprule
		\textbf{Jel} & \textbf{Jelentés} & \textbf{Egység} \\
		\midrule
		\endhead
		$\mbody$ & a főtest tömege & \si{kg} \\ \hline
		$\mwheel$ & egy kerék tömege & \si{kg} \\ \hline
		$r$ & keréksugár & \si{m} \\ \hline
		$b$ & nyomtáv, vagyis a két kerék érintkezési síkja közötti távolság & \si{m} \\ \hline
		$l$ & a keréktengely és a főtest tömegközéppontja közötti távolság & \si{m} \\ \hline
		$I_x$ & a főtest tehetetlenségi nyomatéka a test $x$ tengelye körül & \si{kg.m^2} \\ \hline
		$I_y$ & a főtest tehetetlenségi nyomatéka a test $y$ tengelye körül & \si{kg.m^2} \\ \hline
		$I_z$ & a főtest tehetetlenségi nyomatéka a test $z$ tengelye körül & \si{kg.m^2} \\ \hline
		$\Jspin$ & egy kerék tengelyirányú forgási tehetetlensége & \si{kg.m^2} \\ \hline
		$\Jzw$ & egy kerék elfordulási mozgáshoz tartozó tehetetlensége & \si{kg.m^2} \\ \hline
		$g$ & gravitációs gyorsulás & \si{m.s^{-2}} \\ \hline
		$\tauL$ & bal kerék kerékoldali hajtónyomatéka & \si{N.m} \\ \hline
		$\tauR$ & jobb kerék kerékoldali hajtónyomatéka & \si{N.m} \\
		\bottomrule
	\end{longtable}
	
	\subsubsection{A redukált koordináta-rendszer}
	
	\begin{equation}q_r = \begin{bmatrix} s \\ \theta \\ \psi \end{bmatrix}
	\end{equation}
	
	Így a sebességvektor:
	
	\begin{equation}
		\dot{q_r} = \begin{bmatrix} \dot{s} \\ \dot{\theta } \\ \dot{\psi } \end{bmatrix}
		= \begin{bmatrix} v \\ \dot{\theta } \\ \dot{\psi } \end{bmatrix}
	\end{equation}
	
	
	\begin{itemize}
		\item $v$: a robot előremeneteli sebessége;
		\item $\dot{\theta}$: a főtest dőlési szögsebessége;
		\item $\dot{\psi}$: a robot fordulási szögsebessége 
	\end{itemize}
	
	\newpage
	
	\subsubsection{A keréktengely-középpont síkbeli kinematikája}
	
	\begin{equation}
		e_f = \begin{bmatrix}
			cos(\psi) \\ sin(\psi)
		\end{bmatrix}, \quad
		v_p =
		\begin{bmatrix}
			\dot{x} \\ \dot{y}
		\end{bmatrix}
	\end{equation}
	
	
	
	Oldalirányú csúszás nélkül a sebességvektor a robot előreirányába mutat, ezért:
	
	\begin{equation}
		\text{v}_P = v e_f
	\end{equation}
	
	
	Ebből következik, hogy:
	
	\begin{equation}
	\dot{x} = v \cdot cos(\psi), \quad \dot{y} = v \cdot sin(\psi)
	\end{equation}
	
	\subsubsection{A két kerék gördülési és fordulási kinematikája}

	
	\begin{equation}
		v_L = v - \frac{b}{2} \dot{\psi}, \quad v_R = v - \frac{b}{2} \dot{\psi}
	\end{equation}
	
	Csúszásmentes gördülés esetén \(v_i=r\omega_i=r\dot\phi_i\), ezért:
	
	\begin{equation}
		r\dot\phi_L=v-\frac{b}{2}\dot\psi,
		\qquad
		r\dot\phi_R=v+\frac{b}{2}\dot\psi
	\end{equation}
	
	
	A két egyenlet összeadásából és kivonásából:
	
	\begin{equation}
		\boxed{v=\frac{r}{2}\left(\dot\phi_R+\dot\phi_L\right)},
		\qquad
		\boxed{\dot\psi=\frac{r}{b}\left(\dot\phi_R-\dot\phi_L\right)}.
	\end{equation}
	
	\subsubsection{Az abszolút és relatív kerékszög kapcsolata}
	A relatív encoder-szög a kerék és a robot főtestének szöghelyzete közötti különbség:
	\begin{equation}
	\alpha_i=\phi_i-\theta.
	\end{equation}
	Idő szerint deriválva:
	
	\begin{equation}
	\dot\alpha_i=\dot\phi_i-\dot\theta.	
	\end{equation}
	Mivel \(\omega_i=\dot\phi_i\), ezért:
	
	\begin{equation}
		\boxed{\dot\alpha_L=\omega_L-\dot\theta,\qquad
			\dot\alpha_R=\omega_R-\dot\theta.}
	\end{equation}


\end{document}
